\chapter{Fazit und Ausblick}
\label{sec:fazit_ausblick}

\section{Zusammenfassung der Arbeit}
\label{subsec:zusammenfassung}
% Inhalt: Ein kurzer, prägnanter Rückblick auf die Ausgangslage. Was war das initiale Problem? (Zeit- und fehlerintensive manuelle Migration von SAP LO-VC nach AVC).
% Lösungsansatz: Kurze Zusammenfassung der Methodik (Entwicklung eines neuro-symbolischen Agentensystems auf Basis von LLMs).
% Kernergebnisse: Zusammenfassung der Evaluierungsergebnisse. Betonung der Stärken (hervorragende Extraktion von Logik und unklassifizierten Merkmalen) sowie der aktuellen Schwächen in der Tabellenarchitektur.

\section{Beantwortung der Forschungsfragen}
\label{subsec:beantwortung_rqs}
% Inhalt: Systematische Beantwortung der in Kapitel 1 aufgestellten Forschungsfragen (Research Questions).
% Struktur: Jede RQ (z. B. RQ1, RQ2, etc.) wird idealerweise kurz wiederholt und in 1-2 kompakten Sätzen basierend auf den gewonnenen Erkenntnissen der Arbeit präzise beantwortet.

\section{Zukünftige Entwicklungen (Future Work)}
\label{subsec:future_work}
% Inhalt: Ausblick auf die notwendigen nächsten Schritte zur Weiterentwicklung des Artefakts.
% Kurzfristige technische Optimierungen: Integration des Wildcard-Konzepts (Platzhalter), konsequentes Löschen unklassifizierter Merkmale aus den Tabellendefinitionen und Verschärfung der syntaktischen Prompt-Restriktionen (z. B. Vermeidung von NOT SPECIFIED).
% Langfristige methodische Weiterentwicklung: Ausblick auf eine quantitative Evaluierung der kommenden Version (V2). Notwendigkeit einer "Ground Truth" (ein bereits manuell migriertes Projekt) zur Messung der exakten prozentualen Zeitersparnis und Fehlerquote unter realen Produktionsbedingungen.